\chapter{2009年江苏卷}

\section{填空题}

\begin{problem}
若复数 \(z_1=4+29\i\)，\(z_2=6+9\i\)，其中 \(\i\) 是虚数单位，则复数 \((z_1-z_2)\i\) 的实部为\fillinblank{}.
\end{problem}

\begin{answer}
\(-20\).
\end{answer}

\begin{problem}
已知向量 \(\bs{a}\) 和向量 \(\bs{b}\) 的夹角为 \(30^\circ\)，\(|\bs{a}|=2\)，\(|\bs{b}|=\sqrt3\)，则向量 \(\bs{a}\) 和向量 \(\bs{b}\) 的数量积 \(\bs{a}\cdot\bs{b}=\fillinblank{}\).
\end{problem}

\begin{answer}
\(3\).
\end{answer}

\begin{problem}
函数 \(f(x)=x^3-15x^2-33x+6\) 的单调减区间为\fillinblank{}.
\end{problem}

\begin{answer}
\((-1,11)\).
\end{answer}

\begin{problem}
函数 \(y=A\sin(\omega x+\varphi)\)（\(A\)，\(\omega\)，\(\varphi\) 为常数，\(A>0\)，\(\omega>0\)）在闭区间 \([-\pi,0]\) 上的图象如图所示，则 \(\omega=\fillinblank{}\).

\bitmapfigure[max width=0.36\linewidth]{img/2009/jiangsu/q4_fig1.png}
\end{problem}

\begin{answer}
\(3\).
\end{answer}

\begin{problem}
现有 \(5\) 根竹竿，它们的长度（单位：\(\text{m}\)）分别为 \(2.5\)，\(2.6\)，\(2.7\)，\(2.8\)，\(2.9\). 若从中一次随机抽取 \(2\) 根竹竿，则它们的长度恰好相差 \(0.3\text{ m}\) 的概率为\fillinblank{}.
\end{problem}

\begin{answer}
\(0.2\).
\end{answer}

\begin{problem}
某校甲，乙两个班级各有 \(5\) 名编号为 \(1\)，\(2\)，\(3\)，\(4\)，\(5\) 的学生进行投篮练习，每人投 \(10\) 次，投中的次数如下表：

\begin{center}
\begin{tabular}{c|ccccc}
学生 & \(1\)号 & \(2\)号 & \(3\)号 & \(4\)号 & \(5\)号 \\
\hline
甲班 & \(6\) & \(7\) & \(7\) & \(8\) & \(7\) \\
乙班 & \(6\) & \(7\) & \(6\) & \(7\) & \(9\)
\end{tabular}
\end{center}

则以上两组数据的方差中较小的一个为 \(s^2=\fillinblank{}\).
\end{problem}

\begin{answer}
\(0.4\).
\end{answer}

\begin{problem}
如图是一个算法的流程图，最后输出的 \(W=\fillinblank{}\).

\texfigure[max width=0.55\linewidth]{img_repaint/2009/jiangsu/q7_flow.tex}
\end{problem}

\begin{answer}
\(22\).
\end{answer}

\begin{problem}
在平面上，若两个正三角形的边长的比为 \(1:2\)，则它们的面积比为 \(1:4\)；类似地，在空间内，若两个正四面体的棱长的比为 \(1:2\)，则它们的体积比为\fillinblank{}.
\end{problem}

\begin{answer}
\(1:8\).
\end{answer}

\begin{problem}
在平面直角坐标系 \(xOy\) 中，点 \(P\) 在曲线 \(C:y=x^3-10x+3\) 上，且在第二象限内，已知曲线 \(C\) 在点 \(P\) 处的切线的斜率为 \(2\)，则点 \(P\) 的坐标为\fillinblank{}.
\end{problem}

\begin{answer}
\((-2,15)\).
\end{answer}

\begin{problem}
已知 \(a=\frac{\sqrt5-1}{2}\)，函数 \(f(x)=a^x\)，若实数 \(m\)，\(n\) 满足 \(f(m)>f(n)\)，则 \(m\)，\(n\) 的大小关系为\fillinblank{}.
\end{problem}

\begin{answer}
\(m<n\).
\end{answer}

\begin{problem}
已知集合 \(A=\{x\mid \log_2 x\le2\}\)，\(B=(-\infty,a)\)，若 \(A\subseteq B\)，则实数 \(a\) 的取值范围是 \((c,+\infty)\)，其中 \(c=\fillinblank{}\).
\end{problem}

\begin{answer}
\(4\).
\end{answer}

\begin{problem}
设 \(\alpha\) 和 \(\beta\) 为不重合的两个平面，给出下列命题：

\begin{circlelist}
\item 若 \(\alpha\) 内的两条相交直线分别平行于 \(\beta\) 内的两条直线，则 \(\alpha\) 平行于 \(\beta\)；
\item 若 \(\alpha\) 外一条直线 \(l\) 与 \(\alpha\) 内的一条直线平行，则 \(l\) 和 \(\alpha\) 平行；
\item 设 \(\alpha\) 和 \(\beta\) 相交于直线 \(l\)，若 \(\alpha\) 内有一条直线垂直于 \(l\)，则 \(\alpha\) 和 \(\beta\) 垂直；
\item 直线 \(l\) 与 \(\alpha\) 垂直的充分必要条件是 \(l\) 与 \(\alpha\) 内的两条直线垂直.
\end{circlelist}

上面命题中，真命题的序号是\fillinblank{}.（写出所有真命题的序号）
\end{problem}

\begin{answer}
\(\circled{1}\circled{2}\).
\end{answer}

\begin{problem}
如图，在平面直角坐标系 \(xOy\) 中，\(A_1\)，\(A_2\)，\(B_1\)，\(B_2\) 为椭圆 \(\frac{x^2}{a^2}+\frac{y^2}{b^2}=1\) （\(a>b>0\)）的四个顶点，\(F\) 为其右焦点. 直线 \(A_1B_2\) 与直线 \(B_1F\) 相交于点 \(T\)，线段 \(OT\) 与椭圆的交点 \(M\) 恰为线段 \(OT\) 的中点，则该椭圆的离心率为\fillinblank{}.

\bitmapfigure[max width=0.52\linewidth]{img/2009/jiangsu/q13_fig1.png}
\end{problem}

\begin{answer}
\(2\sqrt{7}-5\).
\end{answer}

\begin{problem}
设 \(\{a_n\}\) 是公比为 \(q\) 的等比数列，\(|q|>1\). 令 \(b_n=a_n+1\)（\(n=1,2,\cdots\)），若数列 \(\{b_n\}\) 有连续四项在集合 \(\{-53,-23,19,37,82\}\) 中，则 \(6q=\fillinblank{}\).
\end{problem}

\begin{answer}
\(-9\).
\end{answer}

\section{解答题}

\begin{problem}
设向量 \(\bs{a}=(4\cos\alpha,\sin\alpha)\)，\(\bs{b}=(\sin\beta,4\cos\beta)\)，\(\bs{c}=(\cos\beta,-4\sin\beta)\).

\begin{enumerate}
\item 若 \(\bs{a}\) 与 \(\bs{b}-2\bs{c}\) 垂直，求 \(\tan(\alpha+\beta)\) 的值；
\item 求 \(|\bs{b}+\bs{c}|\) 的最大值；
\item 若 \(\tan\alpha\tan\beta=16\)，求证：\(\bs{a}\parallel\bs{b}\).
\end{enumerate}
\end{problem}

\begin{solution}
\begin{enumerate}
\item 由垂直条件得
\[
0=\bs{a}\cdot(\bs{b}-2\bs{c})
=4\cos\alpha(\sin\beta-2\cos\beta)
+\sin\alpha(4\cos\beta+8\sin\beta).
\]
整理得
\[
0=4\bigl(\sin(\alpha+\beta)-2\cos(\alpha+\beta)\bigr).
\]
若 \(\cos(\alpha+\beta)=0\)，上式不可能成立，所以
\[
\tan(\alpha+\beta)=2.
\]

\item
\[
|\bs{b}+\bs{c}|^2
=(\sin\beta+\cos\beta)^2+4^2(\cos\beta-\sin\beta)^2
=17-15\sin2\beta\le32.
\]
当 \(\sin2\beta=-1\) 时等号成立，故 \(|\bs{b}+\bs{c}|\) 的最大值为 \(4\sqrt2\).

\item 因为 \(\tan\alpha\) 和 \(\tan\beta\) 有意义，所以 \(\cos\alpha\ne0\)，\(\cos\beta\ne0\). 又
\[
\det(\bs{a},\bs{b})
=4\cos\alpha\cdot4\cos\beta-\sin\alpha\sin\beta
=\cos\alpha\cos\beta\bigl(16-\tan\alpha\tan\beta\bigr)=0.
\]
因此 \(\bs{a}\parallel\bs{b}\).
\end{enumerate}
\end{solution}

\begin{problem}
如图，在直三棱柱 \(ABC-A_1B_1C_1\) 中，\(E\)，\(F\) 分别是 \(A_1B\)，\(A_1C\) 的中点，点 \(D\) 在 \(B_1C_1\) 上，\(A_1D\perp B_1C\). 求证：

\begin{enumerate}
\item \(EF\parallel\) 平面 \(ABC\)；
\item 平面 \(A_1FD\perp\) 平面 \(BB_1C_1C\).
\end{enumerate}

\bitmapfigure[max width=0.42\linewidth]{img/2009/jiangsu/q16_fig1.png}
\end{problem}

\begin{solution}
\begin{enumerate}
\item 在三角形 \(A_1BC\) 中，\(E\)，\(F\) 分别是 \(A_1B\)，\(A_1C\) 的中点，所以由中位线定理得 \(EF\parallel BC\). 又 \(BC\subset\) 平面 \(ABC\)，且 \(EF\) 不在平面 \(ABC\) 内，故 \(EF\parallel\) 平面 \(ABC\).

\item 因为 \(A_1D\) 在平面 \(A_1B_1C_1\) 内，而平面 \(A_1B_1C_1\parallel\) 平面 \(ABC\)，又 \(BB_1\perp\) 平面 \(ABC\)，所以 \(A_1D\perp BB_1\). 题设给出 \(A_1D\perp B_1C\). 直线 \(BB_1\)，\(B_1C\) 是平面 \(BB_1C_1C\) 内的两条相交直线，因此 \(A_1D\perp\) 平面 \(BB_1C_1C\). 又 \(A_1D\subset\) 平面 \(A_1FD\)，所以平面 \(A_1FD\perp\) 平面 \(BB_1C_1C\).
\end{enumerate}
\end{solution}

\begin{problem}
设 \(\{a_n\}\) 是公差不为零的等差数列，\(S_n\) 为其前 \(n\) 项和，满足 \(a_2^2+a_3^2=a_4^2+a_5^2\)，\(S_7=7\).

\begin{enumerate}
\item 求数列 \(\{a_n\}\) 的通项公式及前 \(n\) 项和 \(S_n\)；
\item 试求所有的正整数 \(m\)，使得 \(\frac{a_ma_{m+1}}{a_{m+2}}\) 为数列 \(\{a_n\}\) 中的项.
\end{enumerate}
\end{problem}

\begin{solution}
\begin{enumerate}
\item 设等差数列的公差为 \(d\)，则 \(d\ne0\). 记 \(a_3=u\)，则
\(a_2=u-d\)，\(a_4=u+d\)，\(a_5=u+2d\).
由题设
\[
(u-d)^2+u^2=(u+d)^2+(u+2d)^2,
\]
化简得 \(-4d(2u+d)=0\). 因为 \(d\ne0\)，所以 \(2a_3+d=0\)，即 \(2a_1+5d=0\). 另一方面，
\[
S_7=\frac{7}{2}(2a_1+6d)=7(a_1+3d)=7,
\]
所以 \(a_1+3d=1\). 解得 \(d=2\)，\(a_1=-5\). 因此
\(a_n=a_1+(n-1)d=2n-7\)，\(S_n=\frac{n}{2}\bigl(2a_1+(n-1)d\bigr)=n^2-6n\).

\item 由上问，
\[
\frac{a_ma_{m+1}}{a_{m+2}}=\frac{(2m-7)(2m-5)}{2m-3}.
\]
令 \(t=2m-3\)，则 \(t\) 是奇数且 \(t\ge-1\)，并且
\[
\frac{a_ma_{m+1}}{a_{m+2}}=\frac{(t-4)(t-2)}{t}=t+\frac{8}{t}-6.
\]
若该数是数列中的项，则它必须是整数，所以 \(t\mid8\). 由于 \(t\) 是奇数，只可能有 \(t=1\) 或 \(t=-1\). 当 \(t=1\) 时，\(m=2\)，上式的值为 \(3=a_5\)；当 \(t=-1\) 时，\(m=1\)，上式的值为 \(-15\)，而数列的首项为 \(a_1=-5\)，不可能是数列中的项. 故所有满足条件的正整数为 \(m=2\).
\end{enumerate}
\end{solution}

\begin{problem}
在平面直角坐标系 \(xOy\) 中，已知圆 \(C_1:(x+3)^2+(y-1)^2=4\) 和圆 \(C_2:(x-4)^2+(y-5)^2=4\).

\begin{enumerate}
\item 若直线 \(l\) 过点 \(A(4,0)\)，且被圆 \(C_1\) 截得的弦长为 \(2\sqrt3\)，求直线 \(l\) 的方程；
\item 设 \(P\) 为平面上的点，满足：存在过点 \(P\) 的无穷多对互相垂直的直线 \(l_1\) 和 \(l_2\)，它们分别与圆 \(C_1\) 和圆 \(C_2\) 相交，且直线 \(l_1\) 被圆 \(C_1\) 截得的弦长与直线 \(l_2\) 被圆 \(C_2\) 截得的弦长相等，试求所有满足条件的点 \(P\) 的坐标.
\end{enumerate}

\bitmapfigure[max width=0.42\linewidth]{img/2009/jiangsu/q18_fig1.png}
\end{problem}

\begin{solution}
\begin{enumerate}
\item 设直线 \(l\) 的方程为 \(y=k(x-4)\)，即 \(kx-y-4k=0\). 圆 \(C_1\) 的半径为 \(2\)，弦长为 \(2\sqrt3\)，所以圆心 \((-3,1)\) 到直线 \(l\) 的距离为
\[
\sqrt{2^2-\left(\frac{2\sqrt3}{2}\right)^2}=1.
\]
于是
\[
\frac{|-3k-1-4k|}{\sqrt{k^2+1}}=1,
\]
即 \(24k^2+7k=0\). 因此 \(k=0\) 或 \(k=-\frac{7}{24}\)，所求直线为 \(y=0\) 或 \(7x+24y-28=0\).

\item 设 \(P=(m,n)\). 对于非水平且非竖直的直线，设过 \(P\) 的直线 \(l_1\) 斜率为 \(k\)，则与之垂直的直线 \(l_2\) 斜率为 \(-\frac1k\)，它们的方程可写为
\[
l_1:\ y-n=k(x-m).
\]
\[
l_2:\ y-n=-\frac1k(x-m).
\]
若存在无穷多对符合条件的直线，则符合条件的非零有限斜率 \(k\) 也有无穷多个，因为水平或竖直方向各只有一条直线.

两圆半径相等，且两条直线截得的弦长相等，所以两圆心到相应直线的距离相等. 将 \(C_1=(-3,1)\)，\(C_2=(4,5)\) 代入点到直线的距离公式，得
\[
\frac{|-(m+3)k+n-1|}{\sqrt{k^2+1}}
=\frac{|(n-5)k+m-4|}{\sqrt{k^2+1}}.
\]
因此
\[
(2-m-n)k=m-n-3
\]
或
\[
(m-n+8)k=m+n-5.
\]
关于 \(k\) 的方程只有在相应等式的两边恒为零时才有无穷多个解，所以
\[
\begin{cases}
2-m-n=0,\\
m-n-3=0,
\end{cases}
\quad\text{或}\quad
\begin{cases}
m-n+8=0,\\
m+n-5=0.
\end{cases}
\]
解得 \((m,n)=\left(\frac52,-\frac12\right)\) 或 \((m,n)=\left(-\frac32,\frac{13}{2}\right)\).
对于这两个点，满足距离相等的方向中有无穷多个方向使公共距离不超过半径 \(2\)，故确实能分别与两圆相交. 所以所求点为 \(\left(-\frac32,\frac{13}{2}\right)\) 或 \(\left(\frac52,-\frac12\right)\).
\end{enumerate}
\end{solution}

\begin{problem}
按照某学者的理论，假设一个人生产某产品单件成本为 \(a\text{ 元}\)，如果他卖出该产品的单价为 \(m\text{ 元}\)，则他的满意度为 \(\frac{m}{m+a}\)；如果他买进该产品的单价为 \(n\text{ 元}\)，则他的满意度为 \(\frac{a}{n+a}\). 如果一个人对两种交易（卖出或买进）的满意度分别为 \(h_1\) 和 \(h_2\)，则他对这两种交易的综合满意度为 \(\sqrt{h_1h_2}\). 现假设甲生产 A，B 两种产品的单件成本分别为 \(12\text{ 元}\) 和 \(5\text{ 元}\)，乙生产 A，B 两种产品的单件成本分别为 \(3\text{ 元}\) 和 \(20\text{ 元}\)，设产品 A，B 的单价分别为 \(m_{\text{A}}\text{ 元}\) 和 \(m_{\text{B}}\text{ 元}\)，甲买进 A 与卖出 B 的综合满意度为 \(h_{\text{甲}}\)，乙卖出 A 与买进 B 的综合满意度为 \(h_{\text{乙}}\).

\begin{enumerate}
\item 求 \(h_{\text{甲}}\) 和 \(h_{\text{乙}}\) 关于 \(m_{\text{A}}\)，\(m_{\text{B}}\) 的表达式；当 \(m_{\text{A}}=\frac35m_{\text{B}}\) 时，求证：\(h_{\text{甲}}=h_{\text{乙}}\)；
\item 设 \(m_{\text{A}}=\frac35m_{\text{B}}\)，当 \(m_{\text{A}}\)，\(m_{\text{B}}\) 分别为多少时，甲，乙两人的综合满意度均最大？最大的综合满意度为多少？
\item 记第 \(2\) 问中最大的综合满意度为 \(h_0\)，试问能否适当选取 \(m_{\text{A}}\)，\(m_{\text{B}}\) 的值，使得 \(h_{\text{甲}}\ge h_0\) 和 \(h_{\text{乙}}\ge h_0\) 同时成立，但等号不同时成立？试说明理由.
\end{enumerate}
\end{problem}

\begin{solution}
\begin{enumerate}
\item 甲买进 \(A\) 的满意度为 \(\dfrac{12}{m_{\text{A}}+12}\)，卖出 \(B\) 的满意度为 \(\dfrac{m_{\text{B}}}{m_{\text{B}}+5}\)，所以
\[
h_{\text{甲}}=\sqrt{\frac{12m_{\text{B}}}{(m_{\text{A}}+12)(m_{\text{B}}+5)}}.
\]
乙卖出 \(A\) 的满意度为 \(\dfrac{m_{\text{A}}}{m_{\text{A}}+3}\)，买进 \(B\) 的满意度为 \(\dfrac{20}{m_{\text{B}}+20}\)，所以
\[
h_{\text{乙}}=\sqrt{\frac{20m_{\text{A}}}{(m_{\text{A}}+3)(m_{\text{B}}+20)}}.
\]
当 \(m_{\text{A}}=\frac35m_{\text{B}}\) 时，令 \(m_{\text{B}}=x>0\)，则
\[
h_{\text{甲}}^2=\frac{12x}{\left(\frac35x+12\right)(x+5)}
=\frac{20x}{(x+20)(x+5)},
\]
\[
h_{\text{乙}}^2=\frac{20\cdot\frac35x}{\left(\frac35x+3\right)(x+20)}
=\frac{20x}{(x+20)(x+5)}.
\]
由于两者均为正数，故 \(h_{\text{甲}}=h_{\text{乙}}\).

\item 仍令 \(m_{\text{B}}=x>0\). 由上面的化简，
\[
h_{\text{甲}}^2=h_{\text{乙}}^2=\frac{20x}{(x+20)(x+5)}.
\]
又
\[
(x+20)(x+5)-45x=(x-10)^2\ge0,
\]
所以
\[
h_{\text{甲}}^2=h_{\text{乙}}^2\le\frac49.
\]
等号当且仅当 \(x=10\). 因而当 \(m_{\text{B}}=10\)，\(m_{\text{A}}=6\) 时，甲、乙两人的综合满意度同时取得最大值，且最大值为 \(\frac23\).

\item 由上问 \(h_0=\frac23\). 对任意正的 \(m_{\text{A}}\)，\(m_{\text{B}}\)，有
\[
h_{\text{甲}}^2h_{\text{乙}}^2
=\frac{12m_{\text{A}}}{(m_{\text{A}}+12)(m_{\text{A}}+3)}
\cdot
\frac{20m_{\text{B}}}{(m_{\text{B}}+5)(m_{\text{B}}+20)}.
\]
而
\[
\frac{12m_{\text{A}}}{(m_{\text{A}}+12)(m_{\text{A}}+3)}\le\frac49.
\]
以及
\[
\frac{20m_{\text{B}}}{(m_{\text{B}}+5)(m_{\text{B}}+20)}\le\frac49.
\]
前一个不等式等价于 \((m_{\text{A}}-6)^2\ge0\)，后一个不等式等价于 \((m_{\text{B}}-10)^2\ge0\).
于是 \(h_{\text{甲}}h_{\text{乙}}\le\frac49=h_0^2\). 若同时有 \(h_{\text{甲}}\ge h_0\) 和 \(h_{\text{乙}}\ge h_0\)，则必有 \(h_{\text{甲}}h_{\text{乙}}\ge h_0^2\)，从而两边均取等号. 这要求 \(m_{\text{A}}=6\)，\(m_{\text{B}}=10\)，此时 \(h_{\text{甲}}=h_{\text{乙}}=h_0\). 因此不存在使两个不等式同时成立而等号不同时成立的取值.
\end{enumerate}
\end{solution}

\begin{problem}
设 \(a\) 为实数，函数 \(f(x)=2x^2+(x-a)|x-a|\).

\begin{enumerate}
\item 若 \(f(0)\ge1\)，求 \(a\) 的取值范围；
\item 求 \(f(x)\) 的最小值；
\item 设函数 \(h(x)=f(x)\)，\(x\in(a,+\infty)\)，直接写出（不需给出演算步骤）不等式 \(h(x)\ge1\) 的解集.
\end{enumerate}
\end{problem}

\begin{solution}
\begin{enumerate}
\item 因为
\[
f(0)=-a|a|.
\]
当 \(a\ge0\) 时，\(-a^2\ge1\) 不可能；当 \(a<0\) 时，\(a^2\ge1\)，所以 \(a\le-1\).

\item 当 \(x\ge a\) 时，
\[
f(x)=3x^2-2ax+a^2,
\]
其顶点横坐标为 \(x=\frac a3\). 当 \(a\ge0\) 时，顶点不在区间 \([a,+\infty)\) 内，最小值在 \(x=a\) 处取得，为 \(2a^2\)；当 \(a<0\) 时，顶点在区间 \([a,+\infty)\) 内，最小值为 \(\frac{2a^2}{3}\).

当 \(x<a\) 时，
\[
f(x)=x^2+2ax-a^2.
\]
当 \(a\ge0\) 时，顶点 \(x=-a\) 在该分支的定义域内，最小值为 \(-2a^2\)；当 \(a<0\) 时，该分支在 \(x<a\) 上递减，其下确界为在 \(x=a\) 处的函数值 \(2a^2\)，而右侧分支可在 \(x=a/3\) 处取得更小的值 \(2a^2/3\). 因此
\[
\min_{x\in\R}f(x)=
\begin{cases}
-2a^2,&a\ge0,\\[2pt]
\dfrac{2a^2}{3},&a<0.
\end{cases}
\]

\item 因为 \(x\in(a,+\infty)\)，所以
\[
h(x)=3x^2-2ax+a^2.
\]
不等式 \(h(x)\ge1\) 等价于
\[
3x^2-2ax+a^2-1\ge0.
\]
其判别式为 \(12-8a^2\)，根（存在时）分别为 \(r_-=\frac{a-\sqrt{3-2a^2}}{3}\) 和 \(r_+=\frac{a+\sqrt{3-2a^2}}{3}\).
与定义域 \((a,+\infty)\) 结合，解集为
\[
\{x\in(a,+\infty)\mid h(x)\ge1\}=
\begin{cases}
(a,+\infty),&a\in\left(-\infty,-\dfrac{\sqrt6}{2}\right]\cup\left[\dfrac{\sqrt6}{2},+\infty\right),\\[6pt]
\left[r_+,+\infty\right),&a\in\left[-\dfrac{\sqrt2}{2},\dfrac{\sqrt2}{2}\right),\\[6pt]
\left(a,r_-\right]\cup\left[r_+,+\infty\right),&a\in\left(-\dfrac{\sqrt6}{2},-\dfrac{\sqrt2}{2}\right),\\[6pt]
(a,+\infty),&a\in\left[\dfrac{\sqrt2}{2},\dfrac{\sqrt6}{2}\right).
\end{cases}
\]
\end{enumerate}
\end{solution}

\section{附加题}

\begin{problem}[A]
如图，在四边形 \(ABCD\) 中，\(\triangle ABC\cong\triangle BAD\). 求证：\(AB\parallel CD\).

\bitmapfigure[max width=0.40\linewidth]{img/2009/jiangsu/q21_a_fig1.png}
\end{problem}

\begin{solution}
由 \(\triangle ABC\cong\triangle BAD\)，得 \(\angle ACB=\angle BDA\)，\(\angle CAB=\angle DBA\).
因此点 \(A\)，\(B\)，\(C\)，\(D\) 共圆. 在该圆中，同弦 \(CB\) 所对的圆周角相等，所以
\[
\angle CAB=\angle CDB.
\]
于是 \(\angle DBA=\angle CDB\)，这是一对内错角相等，故 \(AB\parallel CD\).
\end{solution}

\reuseproblemnumber
\begin{problem}[B]
求矩阵 \(A=\begin{bmatrix}3&2\\2&1\end{bmatrix}\) 的逆矩阵.
\end{problem}

\begin{solution}
设 \(A^{-1}=\begin{bmatrix}p&q\\r&s\end{bmatrix}\). 因为
\[
A\begin{bmatrix}p&q\\r&s\end{bmatrix}
=\begin{bmatrix}3p+2r&3q+2s\\2p+r&2q+s\end{bmatrix}
=\begin{bmatrix}1&0\\0&1\end{bmatrix},
\]
所以
\[
\begin{cases}
3p+2r=1,\\
3q+2s=0,\\
2p+r=0,\\
2q+s=1.
\end{cases}
\]
解得 \(p=-1\)，\(q=2\)，\(r=2\)，\(s=-3\). 因此
\[
A^{-1}=\begin{bmatrix}-1&2\\2&-3\end{bmatrix}.
\]
\end{solution}

\reuseproblemnumber
\begin{problem}[C]
已知曲线 \(C\) 的参数方程为 \(x=\sqrt t-\frac{1}{\sqrt t}\)，\(y=3\left(t+\frac1t\right)\)（\(t\) 为参数，\(t>0\)），求曲线 \(C\) 的普通方程.
\end{problem}

\begin{solution}
由参数方程得
\[
x^2=\left(\sqrt t-\frac1{\sqrt t}\right)^2=t+\frac1t-2.
\]
所以 \(t+\frac1t=x^2+2\)，代入 \(y=3\left(t+\frac1t\right)\)，得
\[
y=3x^2+6,
\]
即曲线 \(C\) 的普通方程为
\[
3x^2-y+6=0.
\]
又因为 \(t>0\) 时 \(\sqrt t-1/\sqrt t\) 的取值范围为全体实数，所以该方程没有额外的取值限制.
\end{solution}

\reuseproblemnumber
\begin{problem}[D]
设 \(a\ge b>0\)，求证：\(3a^3+2b^3\ge3a^2b+2ab^2\).
\end{problem}

\begin{solution}
\[
3a^3+2b^3-3a^2b-2ab^2
=3a^2(a-b)-2b^2(a-b)
=(a-b)(3a^2-2b^2).
\]
因为 \(a\ge b>0\)，所以 \(a-b\ge0\)，且
\[
3a^2-2b^2\ge a^2>0.
\]
故 \((a-b)(3a^2-2b^2)\ge0\)，从而
\[
3a^3+2b^3\ge3a^2b+2ab^2.
\]
\end{solution}

\section{解答题}

\begin{problem}
在平面直角坐标系 \(xOy\) 中，抛物线 \(C\) 的顶点在原点，经过点 \(A(2,2)\)，其焦点 \(F\) 在 \(x\) 轴上.

\begin{enumerate}
\item 求抛物线 \(C\) 的标准方程；
\item 求过点 \(F\)，且与直线 \(OA\) 垂直的直线的方程；
\item 设过点 \(M(m,0)\)（\(m>0\)）的直线交抛物线 \(C\) 于 \(D\)，\(E\) 两点，\(ME=2DM\)，记 \(D\) 和 \(E\) 两点间的距离为 \(f(m)\)，求 \(f(m)\) 关于 \(m\) 的表达式.
\end{enumerate}

\bitmapfigure[max width=0.42\linewidth]{img/2009/jiangsu/q22_fig1.png}
\end{problem}

\begin{solution}
\begin{enumerate}
\item 因为抛物线顶点在原点且焦点在 \(x\) 轴上，设其标准方程为 \(y^2=4px\)，则焦点为 \((p,0)\). 将 \(A(2,2)\) 代入得 \(4=8p\)，所以 \(p=\frac12\)，从而抛物线的标准方程为
\[
y^2=2x.
\]

\item 由 \(y^2=2x\) 可知焦点为 \(F\left(\frac12,0\right)\). 直线 \(OA\) 的斜率为 \(1\)，所以过 \(F\) 且垂直于 \(OA\) 的直线斜率为 \(-1\)，其方程为
\[
y=-\left(x-\frac12\right),
\]
即 \(2x+2y-1=0\).

\item 先考虑非竖直且非水平直线，设其斜率为 \(k\ne0\)，用参数 \(t\) 表示直线上的点：
\[
(x,y)=(m+t,kt).
\]
代入抛物线方程得
\[
k^2t^2=2(m+t),
\]
即
\[
k^2t^2-2t-2m=0.
\]
设对应于 \(D,E\) 的根分别为 \(t_D,t_E\). 由韦达定理，\(t_D+t_E=\frac2{k^2}>0\)，\(t_Dt_E=-\frac{2m}{k^2}<0\).
所以两根异号. 由 \(ME=2DM\)，设 \(|t_D|=r\)，\(|t_E|=2r\)，其中 \(r>0\). 结合两根之和为正，必有 \(t_D=-r\)，\(t_E=2r\). 于是 \(r=\frac2{k^2}\)，\(-2r^2=-\frac{2m}{k^2}\).
由第一式得 \(k^2=2/r\)，代入第二式得 \(r=m/2\)，从而 \(k^2=4/m\). 因此
\[
DE=|t_E-t_D|\sqrt{1+k^2}
=3\frac m2\sqrt{1+\frac4m}
=\frac32\sqrt{m(m+4)}.
\]
水平直线 \(y=0\) 与抛物线只有一个交点，不符合题设；竖直直线 \(x=m\) 与抛物线的两个交点关于 \(M\) 对称，满足 \(ME=DM\)，也不满足题设比例，因此不产生其他情形. 故
\[
f(m)=\frac32\sqrt{m(m+4)}.
\]
\end{enumerate}
\end{solution}

\begin{problem}
对于正整数 \(n\ge2\)，用 \(T_n\) 表示关于 \(x\) 的一元二次方程 \(x^2+2ax+b=0\) 有实数根的有序数组 \((a,b)\) 的组数，其中 \(a\)，\(b\in\{1,2,\cdots,n\}\)（\(a\) 和 \(b\) 可以相等）；对于随机选取的 \(a\)，\(b\in\{1,2,\cdots,n\}\)（\(a\) 和 \(b\) 可以相等），记 \(P_n\) 为关于 \(x\) 的一元二次方程 \(x^2+2ax+b=0\) 有实数根的概率.

\begin{enumerate}
\item 求 \(T_{n^2}\) 和 \(P_{n^2}\)；
\item 求证：对任意正整数 \(n\ge2\)，有 \(P_n>1-\frac{1}{\sqrt n}\).
\end{enumerate}
\end{problem}

\begin{solution}
\begin{enumerate}
\item 方程 \(x^2+2ax+b=0\) 有实数根的充要条件是判别式
\[
\Delta=(2a)^2-4b=4(a^2-b)\ge0,
\]
即 \(b\le a^2\). 计算 \(T_{n^2}\) 时，令 \(a=i\)，\(b=j\)，其中 \(i,j\in\{1,2,\ldots,n^2\}\). 当 \(1\le i\le n\) 时，可取的 \(j\) 有 \(i^2\) 个；当 \(n+1\le i\le n^2\) 时，\(i^2\ge n^2\)，所有 \(n^2\) 个 \(j\) 均可取. 所以
\[
T_{n^2}
=\sum_{i=1}^{n}i^2+(n^2-n)n^2
=\frac{n(n+1)(2n+1)}6+n^4-n^3
=\frac{n(6n^3-4n^2+3n+1)}6.
\]
因此
\[
P_{n^2}=\frac{T_{n^2}}{n^4}
=1-\frac{(n-1)(4n+1)}{6n^3}.
\]

\item 在 \(a,b\in\{1,2,\ldots,n\}\) 时，不满足有实数根的数组必须满足 \(b>a^2\)，这要求 \(a^2<n\)，即 \(a<\sqrt n\). 这样的正整数 \(a\) 少于 \(\sqrt n\) 个，而每个 \(a\) 对应的不合格 \(b\) 数量小于 \(n\). 因此不合格数组总数严格小于 \(n\sqrt n\). 于是
\[
1-P_n
=\frac{\text{不合格数组数}}{n^2}
<\frac{n\sqrt n}{n^2}
=\frac1{\sqrt n},
\]
从而
\[
P_n>1-\frac1{\sqrt n}.
\]
\end{enumerate}
\end{solution}
